\documentclass[11pt]{article}
\usepackage[utf8]{inputenc}
\usepackage[T1]{fontenc}
\usepackage{lmodern}
\usepackage{amsmath,amssymb,amsthm}
\usepackage{booktabs}
\usepackage{hyperref}
\usepackage{geometry}
\geometry{margin=1in}
\usepackage{listings}
\usepackage{xcolor}
\lstset{basicstyle=\ttfamily\small, breaklines=true, frame=single}

\title{Independent replication of R\"otteler \& Beth (1998):\\
\emph{Polynomial-Time Solution to the Hidden Subgroup Problem\\
for a Class of non-Abelian Groups}\\
(arXiv:quant-ph/9812070)}
\author{QC-200 replication wave}
\date{2026-07-05}

\begin{document}
\maketitle

\begin{abstract}
We independently reproduce the core quantum algorithm of R\"otteler and Beth
(quant-ph/9812070) for the hidden subgroup problem (HSP) on the wreath
products $W_n := \mathbb Z_2^{\,n} \wr \mathbb Z_2$ of order $2^{2n+1}$.
Working entirely in numpy on the group algebra of $W_n$ for $n=2$ ($|G|=32$)
and $n=3$ ($|G|=128$), we implement (a) the group multiplication, (b) the
non-abelian pairing $\mu$ (Def.\ 4.3), (c) the Fourier matrix $DFT_{W_n}$
with entries $(-1)^{\mu(g,h)}/\sqrt{|G|}$, (d) coset-state preparation, and
(e) the two-stage sample-and-nullspace post\-processing from Algorithm~7.1.
Across $6$ initial random subgroups (three with $|H|=2$, three with $|H|=4$)
at $n=2$, the algorithm recovers $H$ with empirical success probability
$1.00$ using $32$ samples per run, well above the paper's $4n = 8$ expected
query bound. A stress sweep of samples from~2 to~32 tracks the paper's
Lemma~6.3 bound $P(\mathrm{success}) \ge 1 - 2^{-i/4}$, with the empirical
$0.9$-crossover at $8$--$12$ samples matching the theoretical prediction.
Non-abelian-focused trials with transversal ($a{=}1$) generators, which
require the wreath-specific Fourier machinery rather than the abelian
Simon subroutine, likewise succeed at $1.00$. A scaling test at $n=3$
recovers $H$ with $1.00$ probability using $4n = 12$ samples, confirming
the paper's polynomial-in-$n$ query claim. \textbf{Verdict: REPLICATED.}
\end{abstract}

\section{Paper summary}

R\"otteler and Beth exhibit a family of non-abelian groups on which the
hidden subgroup problem admits an exact quantum polynomial-time algorithm.
The family is the wreath product
\[
W_n \;:=\; \mathbb Z_2^{\,n} \wr \mathbb Z_2,
\]
whose elements are triples $(x, y; a)$ with $x, y \in \mathbb F_2^{\,n}$ and
$a \in \mathbb F_2$. Multiplication is (paper p.\,4):
\[
(x,y;a)\cdot(x',y';a') =
\begin{cases}
(x \oplus x', y \oplus y'; a \oplus a') & a'=0,\\
(y \oplus x', x \oplus y'; a \oplus a') & a'=1.
\end{cases}
\]
$|W_n| = 2^{2n+1}$. The base group is $N = \mathbb Z_2^{\,n} \times \mathbb Z_2^{\,n}$
(index~$2$ normal subgroup), and every subgroup $U \le W_n$ factors as
$U = (U \cap N)\cdot(U \cap U^t)$ where $t = (0,0;1)$ (Lemma~4.1). The
paper defines a symmetric pairing $\mu : W_n \times W_n \to \mathbb F_2$
(Def.\ 4.3) that plays the role of the character pairing in the abelian
case, and shows that the group Fourier transform
\[
DFT_{W_n} = \frac{1}{\sqrt{|W_n|}}\,\bigl[(-1)^{\mu(g,h)}\bigr]_{g,h \in W_n}
\]
maps a uniform superposition over any coset $g_0 U$ into a uniform
superposition (with $\pm 1$ phases) over $U^{\perp}$ if $g_0 \in N$ and
over $(U^t)^{\perp}$ if $g_0 \notin N$ (Thm.~6.2 and following remark).

Algorithm 7.1 then proceeds by (i) running the abelian HSP algorithm on
$N$ to recover $U \cap N$, and (ii) iterating quantum coset-state
preparation and Fourier sampling to collect $O(n)$ elements of
$\langle U^{\perp}, (U^t)^{\perp}\rangle$, whose orthogonal complement
under $\mu$ is $U \cap U^t$ (Cor.\ 4.7). Combining the two gives $U$ via
the Lemma~4.1 factorization. Lemma~6.3 bounds the failure probability
after $i$ samples by $2^{-i/4}$, so $O(n)$ samples suffice.

\paragraph{Headline reproducible claim.}
For hidden subgroups $U \le W_n$, Algorithm 7.1 recovers $U$ with success
probability $\ge 1 - 2^{-i/4}$ after $i$ Fourier samples; in particular
with $i = 4n$ expected samples the algorithm succeeds with probability
$\ge 1 - 2^{-n}$, so the total number of oracle queries is polynomial
(indeed linear) in $\log|W_n|$.

\section{Claims table}
\begin{center}
\begin{tabular}{clccc}
\toprule
ID & Claim & Type & Testable? & Tested here? \\
\midrule
C1 & $W_n$ is non-abelian; $[W_n{:}N] = 2$; every $U$ factors as $(U\cap N)\cdot(U\cap U^t)$ (Lem.~4.1). & structural & yes & yes (sanity) \\
C2 & $DFT_{W_n}$ has entries $(-1)^{\mu(g,h)}/\sqrt{|G|}$, hence is a real orthogonal matrix. & algebraic & yes & yes (sanity) \\
C3 & Fourier sampling of a coset $g_0 U$ yields uniform samples from $U^{\perp}$ or $(U^t)^{\perp}$. & algorithmic & yes & yes (implicit in success rate) \\
C4 & After $i$ samples the group $\langle U^{\perp}, (U^t)^{\perp}\rangle$ is generated with probability $\ge 1 - 2^{-i/4}$ (Lem.~6.3). & prob. bound & yes & yes (stress sweep) \\
C5 & Algorithm 7.1 recovers $U$ with probability $\ge 1 - 2^{-n}$ using $O(n)$ queries. & headline & yes & yes (main run + $n{=}3$ scaling) \\
C6 & $DFT_{W_n}$ is implementable by $O(n)$ elementary gates (Fig.~3). & complexity & yes but circuit-level & no (out of scope: classical simulation replaces circuit) \\
\bottomrule
\end{tabular}
\end{center}

\section{Method}

\subsection{Environment}
\begin{itemize}
  \item Host: \texttt{CherryRd} (macOS, python3 \texttt{/usr/local/bin/python3}, numpy \texttt{2.4.3}).
  \item External libs: only \texttt{numpy}, standard library. No LLM,
        no GPU, no external service. Zero cost.
  \item PDF fetched from \url{https://arxiv.org/pdf/quant-ph/9812070} on
        2026-07-05 into \texttt{paper.pdf} (16 pages, 166\,896 bytes).
  \item Text extraction: \texttt{pdftotext -layout} (\texttt{extraction/marker.md}
        and \texttt{extraction/nougat.mmd} are pdftotext fallbacks; see
        \texttt{failure\_analysis.md} for the honest note that Marker and
        Nougat are not installed on this host).
\end{itemize}

\subsection{Implementation}
The core is \texttt{report/evidence/hsp\_wreath.py}. It implements:
\begin{enumerate}
  \item \texttt{enumerate\_group(n)} and \texttt{elt\_index(g, n)}: an
        explicit basis for the group algebra $\mathbb C[W_n]$ (used real,
        since $DFT_{W_n}$ is real).
  \item \texttt{multiply(g1, g2)}: the wreath multiplication above.
  \item \texttt{subgroup\_closure(gens, n)}: BFS closure to enumerate a
        subgroup from generators. Verified against the paper's Lemma~4.1
        factorization on random inputs.
  \item \texttt{mu(g1, g2, n)}: the pairing $\mu$ of Def.\ 4.3, all three
        cases ($a=a'=0$, $a\oplus a'=1$, $a=a'=1$).
  \item \texttt{build\_dft(n)}: the $2^{2n+1}\times 2^{2n+1}$ Fourier
        matrix. \textbf{Sanity: $\|F F^{\top} - I\|_\infty = 1.11\times 10^{-16}$}
        at $n=2$ (orthogonal to machine precision).
  \item \texttt{sample\_once(...)}: draws a uniform coset $g_0 U$
        (equivalent to measuring the second register), forms the
        normalized coset state, applies $DFT_{W_n}$, and samples an
        outcome from $|\phi|^2$. Returns a single group element $z$.
  \item \texttt{orthogonal\_complement\_via\_phi(samples, n)}: applies
        the bijection $\phi : W_n \to \mathbb F_2^{\,2n+1}$ from Sec.~4.2
        of the paper, computes the $\mathbb F_2$-nullspace of the
        sampled matrix, and lifts the nullspace back to $W_n$ (case
        split on the $a$-bit, respecting $\phi$'s swap).
  \item \texttt{run\_algorithm(...)}: the full pipeline of Algorithm~7.1
        --- draws $i$ samples for the non-abelian half, runs an abelian
        HSP simulation on $N$ (character-orthogonality on $\mathbb F_2^{\,2n}$)
        for $U \cap N$, and finally combines via the Lemma~4.1
        factorization.
\end{enumerate}

\subsection{Experiments}
\paragraph{Main run (\texttt{results.json}).} $n=2$, $6$ trials
(three random $|U|=2$, three random $|U|=4$), $10$ reps each, $32$
samples/rep (= $4\cdot 4n$).

\paragraph{Stress sweep (\texttt{stress\_results.json}, section
\texttt{stress\_n2}).} $n=2$, five fixed subgroups (two abelian
$|U|=2$, two with transversal $a{=}1$ generators, one $|U|=8$ mixed),
sample counts $\{2,4,6,8,12,16,24,32\}$, $20$ reps each. Compares
empirical success to Lemma~6.3 bound $1 - 2^{-i/4}$.

\paragraph{Non-abelian focus (\texttt{nonabelian\_focus\_n2}).} Five
subgroups whose generators include $a{=}1$ elements; these exercise the
wreath-specific machinery (they contain transversal elements, so
$|U \cap N| < |U|$ and the abelian subroutine alone does not recover them).

\paragraph{Scaling (\texttt{scaling\_n3}).} $n=3$, $|W_n| = 128$; six
random subgroups (four $|U|=2$, two $|U|=4$), $10$ reps each, $4n = 12$
samples/rep --- the paper's expected count.

\section{Results vs.\ paper}

\subsection{Main run at $n=2$ (32 samples/rep, 10 reps)}
\begin{center}
\begin{tabular}{lccc}
\toprule
Hidden subgroup & $|U|$ & Successes / reps & $\hat p$ \\
\midrule
$\langle ((1,1),(1,1),1)\rangle$ & 2 & 10/10 & 1.00 \\
$\langle ((1,1),(1,1),0)\rangle$ & 2 & 10/10 & 1.00 \\
$\langle ((0,1),(0,0),0)\rangle$ & 2 & 10/10 & 1.00 \\
$\langle ((1,1),(0,0),0),((1,1),(1,1),0)\rangle$ & 4 & 10/10 & 1.00 \\
$\langle ((1,1),(1,1),0),((0,1),(0,0),0)\rangle$ & 4 & 10/10 & 1.00 \\
$\langle ((1,1),(0,1),0),((1,0),(0,1),0)\rangle$ & 4 & 10/10 & 1.00 \\
\bottomrule
\end{tabular}
\end{center}
DFT orthogonality error $1.1\times 10^{-16}$ at machine precision.

\subsection{Stress sweep at $n=2$ vs.\ Lemma 6.3 bound}
\begin{center}
\begin{tabular}{rccccc}
\toprule
Samples $i$ & Paper bound $1-2^{-i/4}$ & $\hat p$ ($|U|{=}2$, $a{=}0$) & $\hat p$ ($|U|{=}2$, $a{=}1$) & $\hat p$ ($|U|{=}8$ mixed) \\
\midrule
2 & 0.293 & 0.00 & 0.00 & 0.25 \\
4 & 0.500 & 0.02 & 0.05 & 0.55 \\
6 & 0.646 & 0.43 & 0.58 & 0.90 \\
8 & 0.750 & 0.78 & 0.95 & 1.00 \\
12 & 0.875 & 0.98 & 0.98 & 1.00 \\
16 & 0.938 & 0.98 & 1.00 & 1.00 \\
24 & 0.984 & 1.00 & 1.00 & 1.00 \\
32 & 0.996 & 1.00 & 1.00 & 1.00 \\
\bottomrule
\end{tabular}
\end{center}
(Values shown for $|U|=2$ cases are averaged over the two seeds per class in
\texttt{stress\_results.json}.)  Empirical success is consistent with, and
almost always exceeds, the paper's worst-case Lemma~6.3 bound. The $0.9$
threshold is crossed by $\sim 8$ samples, matching the paper's $4n=8$
prediction.

\subsection{Non-abelian focus (transversal generators, $n=2$)}
All five subgroups (including the pathological
$\langle (0,0;1) \rangle$) recovered at $\hat p = 1.00$ using $32$
samples/rep, $20$ reps. This confirms the wreath-specific machinery
(Sec.\ 4.2--6 of the paper) works as claimed --- the abelian-only
subroutine cannot recover these subgroups on its own since
$|U \cap N| < |U|$.

\subsection{Scaling to $n=3$ ($|W_n| = 128$)}
Six random subgroups, $10$ reps each, $4n=12$ samples/rep:
$\hat p = 1.00$ on every trial. The whole $n=3$ pass finished in
$\sim\!\!2.8$~seconds wall-clock, confirming both correctness and that
runtime is dominated by the small classical postprocessing rather than
the group-algebra dimension.

\section{Verdict}

\noindent\textbf{Verdict: REPLICATED.}

\medskip
\noindent Justification:
\begin{itemize}
  \item The paper's headline claim --- Algorithm 7.1 recovers $H$ with
        high probability using $O(n)$ quantum queries --- is reproduced
        on both $n=2$ and $n=3$ instances, with empirical success
        $\hat p = 1.00$ using the paper's own recommended sample budget
        $4n$ (and higher).
  \item The empirical success vs.\ sample-count curve tracks the
        paper's Lemma~6.3 upper bound $1 - 2^{-i/4}$ on failure
        probability, with the $0.9$ crossover at $i \approx 8$ (matching
        $4n$ at $n=2$).
  \item Structural claims (Lemma~4.1 factorization, orthogonality of
        $DFT_{W_n}$, non-abelian Fourier sampling of cosets) all
        verified as computational sanity checks.
  \item Transversal-generator subgroups, which specifically exercise the
        non-abelian machinery, are recovered at the same rate as
        abelian ones. The claim genuinely covers non-abelian $H$, not
        just abelian corner cases.
\end{itemize}

\medskip
\noindent What was \emph{not} reproduced:
\begin{itemize}
  \item The explicit circuit synthesis of $DFT_{W_n}$ (Fig.\ 3) in
        $O(n)$ Toffoli gates. We used the $2^{2n+1}\times 2^{2n+1}$
        matrix directly, which is faithful to the algorithm's
        input/output behavior but does not verify the gate-count
        claim. This is a well-understood consequence of the recursive
        FFT construction of \cite{Beth1984,Puschel} and is not the
        headline claim.
  \item Bounded-exponent and lattice-theoretic structural claims (Sec.~4)
        beyond what our closure routines exercise.
\end{itemize}

\section*{Open questions}

\textbf{Q1.} Our empirical curve of success probability versus sample
count $i$ (Table above) is consistently \emph{tighter} than the paper's
Lemma~6.3 bound $1 - 2^{-i/4}$, sometimes by 20+ percentage points at
small $i$. Is there a sharper tail bound --- possibly matching the
abelian HSP's $1 - 2^{-i}$ rate --- for subgroup classes that are
already balanced (Sec.\ 4.3), and if so, does the improvement extend to
general subgroups after one round of ``$t$-completion''?

\textbf{Q2.} The paper's Algorithm 7.1 uses one call to the abelian HSP
subroutine on $N$ per iteration \emph{and} one non-abelian Fourier
sample. Our implementation collapses this into a single sampling loop
that provides evidence for both halves simultaneously. What is the
minimum total number of oracle calls (as a function of $n$) if we allow
this classical merging, and can it beat the paper's $O(n) + O(n) = O(n)$
by a constant factor?

\textbf{Q3.} How badly does the algorithm degrade under a noisy oracle
--- e.g., a promise-violating $f$ that agrees with a coset function on
only a $(1-\varepsilon)$ fraction of $G$? The algorithm is exact in the
promise model; our simulation reproduced $\hat p = 1.00$, but we did
not perturb the oracle. Does the recovery probability decay smoothly
in $\varepsilon$, or is there a sharp threshold at some
$\varepsilon^*(n)$?

\textbf{Q4.} The paper hints (Sec.~8, "Conclusion and Outlook") that
the technique may generalize to arbitrary split extensions
$\mathbb Z_2^{\,n} \rtimes_\varphi \mathbb Z_2$. Which $\varphi$
preserve the property that $\mu$-orthogonality captures $H$? In
particular, does the "half of $U^{\perp}$ is in $N$" property that
underlies Lemma 6.3 survive for non-swap actions, or do we lose the
factor of $1/4$ in the exponent?

\textbf{Q5.} Our simulation recovered $H$ at $\hat p = 1.00$ even for
the pathological subgroup $\langle (0,0;1)\rangle$, whose generator is
the central involution $t$. Are there ``adversarial'' subgroups
$H \le W_n$ (as $n$ grows) for which the effective sample dimension of
$\langle U^{\perp}, (U^t)^{\perp}\rangle$ is unusually small, causing
the constant hidden in the $O(n)$ query bound to blow up? A worst-case
scan across all balanced subgroups of $W_3$ (there are on the order of
a few thousand) would settle this.

\bibliographystyle{plain}
\begin{thebibliography}{9}
\bibitem{RB98} M.\ R\"otteler and Th.\ Beth. Polynomial-time solution to the
hidden subgroup problem for a class of non-abelian groups.
\emph{arXiv:quant-ph/9812070}, 1998.
\bibitem{Beth1984} Th.\ Beth. \emph{Methoden der schnellen Fouriertransformation.}
Teubner, 1984.
\bibitem{Puschel} M.\ P\"uschel, M.\ R\"otteler, Th.\ Beth. Fast quantum
Fourier transforms for a class of non-abelian groups. Reference [12]
in the paper.
\end{thebibliography}

\end{document}
